%!TEX root = ../dolgozat.tex
\chapter{Bevezetés}
\label{ch:bevezetes}

A felsőoktatásban az elmúlt években alapvetően megváltozott a számonkérés módja. A
COVID-19 járvány idején az egyetemek kényszerből tértek át az online oktatásra és vizsgáztatásra,
a járvány elmúltával azonban sok intézmény megtartotta a digitális számonkérés bevált elemeit. Az
online vizsga gyorsabb értékelést, egységesebb pontozást és rugalmasabb szervezést tesz lehetővé, a
nagy létszámú évfolyamok oktatói számára pedig jelentős tehermentesítést jelent. Ugyanakkor
olyan kérdéseket is felvet, amelyekre a hagyományos, papíralapú vizsgáknál nem kellett választ
keresni: hogyan lehet megbízhatóan, igazságosan és biztonságosan értékelni a hallgatók munkáját,
és hogyan lehet meggyőződni arról, hogy a vizsgázó valóban önállóan dolgozott.

Különösen igaz ez az informatikai képzésekre. Egy programozási feladat megoldását nem elég
,,elolvasni'', a programot le is kell futtatni, és meg kell vizsgálni, hogy különböző bemenetekre
helyes eredményt ad-e. Ha ezt az oktató kézzel végzi, a javítás rendkívül időigényes, és a
pontozás könnyen következetlenné válik: ugyanazt a hibát az egyik dolgozatban észreveszi, a
másikban nem. Kézenfekvő megoldásnak tűnik a programok automatikus, tesztesetekkel történő
ellenőrzése, ez azonban újabb nehézséget hoz magával. A szervernek ilyenkor egy ismeretlen
személy által írt, tetszőleges programkódot kell futtatnia, amely hibából vagy szándékosan
megpróbálhat kárt tenni a rendszerben: végtelen ciklusba kerülhet, elfogyaszthatja a memóriát,
vagy hozzáférhet olyan adatokhoz, amelyekhez nem volna szabad.

Témaválasztásomat elsősorban saját hallgatói tapasztalataim motiválták. A tanulmányaim során
több online vizsgáztató rendszert is használtam, és azt tapasztaltam, hogy ezek vagy általános
célú tesztrendszerek, amelyek a programozási feladatok értékelésére nem alkalmasak, vagy
kifejezetten programozási versenyekre készült megoldások, amelyekből hiányoznak az egyetemi
vizsga megszokott elemei, például a vegyes feladattípusok és a kézi javítás lehetősége. Érdekelt,
hogy egy modern webes technológiákra épülő, átlátható felépítésű rendszerben hogyan lehet ezt a
kettőt összekapcsolni, és hogyan lehet a nem megbízható kód futtatását a gyakorlatban is
biztonságossá tenni.

Dolgozatom célja egy olyan webes vizsgáztató rendszer prototípusának megtervezése és
elkészítése, amely egy vizsgán belül különböző típusú feladatokat kezel (egy- és több
választós, igaz--hamis, rövid szöveges, kifejtős és programozási feladatokat), a programozási
feladatokat izolált Docker-konténerben, nyilvános és rejtett tesztesetekkel automatikusan
értékeli, és valós időben értesíti az oktatót, ha egy hallgató a vizsga közben elhagyja a
vizsgaablakot. A munka során a következő kérdésekre kerestem a választ: Milyen eszközökkel
futtatható biztonságosan egy nem megbízható forrásból származó program? Milyen pontozási
szabályok igazságosak a részben helyes válaszok esetén? Mennyit ér, és hol vannak a határai a
böngészőben megvalósítható vizsgafelügyeletnek?

A dolgozat felépítése a következő. \Aref{ch:koncepcio}. fejezetben a számítógéppel támogatott
értékelés, a programozási feladatok automatikus értékelése és az online vizsgafelügyelet
szakirodalmát tekintem át, bemutatom a hasonló célú rendszereket és a felhasznált
technológiákat, majd megfogalmazom a rendszerrel szemben támasztott követelményeket.
\Aref{ch:tervezes}. fejezet a rendszer architektúráját, adatmodelljét és a főbb folyamatok
tervét ismerteti. \Aref{ch:megvalositas}. fejezetben a megvalósítás érdekesebb részleteit
mutatom be, különös tekintettel a kódfuttatás biztonságára és a valós idejű felügyeletre.
\Aref{ch:teszteles}. fejezet a rendszer funkcionális és biztonsági tesztelésének eredményeit
tartalmazza, végül \aref{ch:osszefoglalas}. fejezetben összegzem az elért eredményeket, és
kitérek a továbbfejlesztési lehetőségekre.
